\subsubsection{Sentiment Analysis}\label{sec:sentiment-analysis}

After seeing how word embeddings improve \textbf{Information Retrieval (IR)} and \textbf{Document Similarity}, another important application is \definition{Sentiment Analysis}. The goal of sentiment analysis is to \hl{automatically determine the emotional polarity expressed by a word, sentence, or document.} For example, ``\emph{I love this move}'' has a \textbf{positive} sentiment, while ``\emph{I hate this movie}'' has a \textbf{negative} sentiment.

\highspace
\begin{flushleft}
    \textcolor{Red2}{\faIcon{exclamation-triangle} \textbf{Traditional Approach}}
\end{flushleft}
Before word embeddings, sentiment analysis was usually formulated as a \textbf{classification problem}. The workflow looked like this:
\begin{enumerate}
    \item Sentence
    \item Bag-of-Words (BoW)/TF-IDF representation
    \item Classifier (e.g., SVM, Logistic Regression, etc.)
    \item Positive/Negative/Neutral label
\end{enumerate}
The classifier had to learn which words were associated with positive or negative emotions. This required labeled training data, feature engineering, and training a supervised model.

\highspace
\begin{flushleft}
    \textcolor{Green3}{\faIcon{check-circle} \textbf{Word Embeddings offer another possibility}}
\end{flushleft}
Because Word2Vec organizes words according to their \textbf{contexts}, words expressing similar emotions naturally become close in the embedding space. For example, ``\emph{happy}'', ``\emph{joyful}'', ``\emph{excited}'', and ``\emph{delighted}'' will tend to occupy the same region. Likewise, ``\emph{depressed}'', ``\emph{sad}'', ``\emph{miserable}'', and ``\emph{heartbreaking}'' will form another cluster. Therefore, instead of training a classifier, we can \hl{simply measure \textbf{how close two words are in the embedding space} to determine their sentiment similarity.}

\highspace
As always, suppose we choose a reference word ``\emph{sad}'', we compute its embedding $\mathbf{w}_{\emph{sad}}$ and compare it with every other word in the vocabulary using the cosine similarity:
\begin{equation*}
    \cos\left(\theta\right) = \frac{\mathbf{w}_{i} \cdot \mathbf{w}_{\emph{sad}}}{\left\|\mathbf{w}_{i}\right\| \left\|\mathbf{w}_{\emph{sad}}\right\|}
\end{equation*}
Words with the highest cosine similarity are considered the most semantically related.

\highspace
\textcolor{Green3}{\faIcon{check-circle} \textbf{Advantages.}} Using word embeddings for sentiment analysis has several advantages:
\begin{itemize}
    \item \textbf{Words with similar emotional meanings} are automatically grouped together;
    \item \textbf{Synonyms} can be \textbf{recognized} even if they never appeared in the labeled training set;
    \item \textbf{Semantic relationships} are captured instead of relying on exact word matching;
    \item \textbf{Cosine similarity} can retrieve emotionally similar words \hl{without explicitly training a sentiment classifier.}
\end{itemize}
\textcolor{Red2}{\faIcon{exclamation-triangle} \textbf{Limitations.}} Although this approach is elegant, it also has important limitations. Word2Vec learns \textbf{word embeddings}, not \textbf{sentence meanings}. Therefore, it cannot correctly handle:
\begin{itemize}
    \item Negation: ``\emph{I am not happy}'' would still be close to ``\emph{happy}'' in the embedding space, leading to incorrect sentiment classification;
    \item Irony or sarcasm: ``\emph{Great, another rainy day}'' might be interpreted as positive due to the presence of ``\emph{Great}'', even though the overall sentiment is negative;
    \item Long-range context: The sentiment of a sentence can depend on the entire context, which word embeddings alone cannot capture.
\end{itemize}
In these cases, modern contextual models such as \textbf{BERT}, \textbf{RoBERTa}, or \textbf{GPT} are much more effective because they analyze the entire sentence rather than individual word embeddings.

\highspace
In summary, Word2Vec embeddings can be used for sentiment analysis because \hl{words expressing similar emotions appear in similar contexts and therefore become neighbors in the embedding space.} By computing cosine similarity between embeddings, we can retrieve emotionally related words without explicitly training a sentiment classifier. This demonstrates once again that Word2Vec learns a \textbf{semantic organization of language}, and that these learned embeddings can be reused across many downstream NLP tasks.